\section{Frobenius 的迹}\srcpage{38}

取有限维向量空间 $V$ 的自同态 $f$ 的迹，这一通常操作有范畴的类比。更确切地说，更一般地，若 $V$ 是某个对称幺半范畴 $\cC$ 中的可对偶对象，且 $f\in\End(V)$，则可以定义
\[ \tr(f|V)\in\End_{\cC}(\mathbf{1}) \]
为复合
\[ \mathbf{1}\to V\otimes V^{*}\xrightarrow{f\otimes\id}V\otimes V^{*}\to\mathbf{1}. \]
这一构造有许多有趣的函子性性质，特别是迹的循环不变性，但我们在此不需要这些性质。

这也可以应用于 $\cC=\PrL_{\e}$ 以及任何可对偶的可表现 $\e$-线性 $(\infty,1)$-范畴。对可表现 $\e$-线性 $(\infty,1)$-范畴而言，可对偶这一条件原来是相当弱但非常有趣的条件；值得注意的是，Efimov（2025）最近把 $K$-理论推广到了这一框架。特别地，所有紧生成范畴都可对偶；并不要求同态空间具有某种有限性。\footnote{反过来，由于 $\PrL_{\e}$ 是 $(\infty,2)$-范畴，还有更强的 $2$-可对偶性概念。$2$-可对偶会推出紧对象之间同态空间的这类有限性条件。} 尽管存在有趣的可对偶非紧生成范畴的例子——例如流形上的层——我们实际上只需要紧生成的情形。

于是，给定任何可对偶的可表现 $\e$-线性 $(\infty,1)$-范畴 $\cC$ 与任何自函子 $F:\cC\to\cC$（在 $\PrL_{\e}$ 中，即保持余极限且 $\e$-线性），我们可以构造
\[ \tr(F|\cC)\in\End_{\PrL_{\e}}(\mathbf{1})=D(\e). \]
换言之，取范畴上的迹，得到的是一个向量空间（或者，由于我们完全在导出框架下工作，是向量空间导出范畴中的一个对象）。

\begin{example}
若 $F=\id$，这恢复 $\cC$ 的 Hochschild 同调。若对某个结合 $\e$-代数 $A$ 有 $\cC=D(A)$，则它由循环 Bar 复形
\[ [\dots\to A\otimes A\otimes A\to A\otimes A\to A] \]
计算，其中微分是 $a\otimes b\mapsto ab-ba$ 与 $a\otimes b\otimes c\mapsto ab\otimes c-a\otimes bc+ca\otimes b$，并对更多因子作适当推广。若 $\cC=D(A)$ 而 $F$ 是一般的，Morita 理论表明 $F$ 由某个 $A$-双模 $M$ 给出，而 $\tr(F|\cC)$ 由
\[ [\dots\to M\otimes A\otimes A\to M\otimes A\to M] \]
计算，其中微分是 $m\otimes a\mapsto ma-am$ 与 $m\otimes a\otimes b\mapsto ma\otimes b-m\otimes ab+ma\otimes b$\yizhu{上式末项原文亦作 $ma\otimes b$，与首项重复；对照前一条公式（末项为循环项 $ca\otimes b$），按循环 Bar 复形的通例此处应为 $mb\otimes a$。译文照录原文。}{原文笔误（第 38 页；上下文推断）}，利用双模结构，并同样作适当推广。
\end{example}

注意，若 $\cC$ 紧生成，则由 Schwede--Shipley 重构定理，它是形如 $D(A)$ 的范畴的递增并。由于迹构造\srcpage{39}与滤余极限交换，这在一般情形给出计算它的手段，至少当 $F$ 保持这一递增并的描述时如此。

可以在（导出的）向量空间 $\tr(F|\cC)$ 中描述一些显式的向量。即，若 $X\in\cC$ 是紧对象并带有一个映射 $\alpha:X\to F(X)$，就得到一个类
\[ [X,\alpha]\in\tr(F|\cC). \]

\begin{theorem}[Arinkin、Gaitsgory、Kazhdan、Raskin、Rozenblyum 与 Varshavsky，2021]
自然映射
\[ \tr(\Frob|D_{\Nilp}(\Bun_G))\to\cA_c(\Bun_G(\Fq)) \]
是同构。
\end{theorem}

利用 $(\infty,2)$-核范畴 $K_D$ 代替 $\PrL_{\e}$，并把迹形式化应用于其中的对象 $D(\Bun_G)$，实际上容易看出
\[ \tr(\Frob|D(\Bun_G))=\cA_c(\Bun_G(\Fq)). \]
事实上，在 $K_D$ 中，对象 $D(\Bun_G)$ 是自对偶的，其自对偶数据由 $\Delta_{!}\Ql\in D(\Bun_G\times\Bun_G)=D(\Bun_G)\otimes_{K_D}D(\Bun_G)$ 与
\[ (M,N)\mapsto\pi_{!}(M\otimes N):D(\Bun_G)\otimes_{K_D}D(\Bun_G)\to\mathbf{1} \]
给出。有了这些显式公式，容易算出 $\tr(\Frob|D(\Bun_G))$ 求值为 $\Bun_G^{\Frob=\id}=\Bun_G(\Fq)$ 的紧支撑上同调，即 $\cA_c(\Bun_G(\Fq))$。

由 $D_{\Nilp}(\Bun_G)$ 用 prim 对象所作的刻画，函子
\[ D_{\Nilp}(\Bun_G)\to D(\Bun_G) \]
在 $K_D$ 中有右伴随，从而诱导迹上的一个映射。这就是被证明为同构的映射
\[ \tr(\Frob|D_{\Nilp}(\Bun_G))\to\cA_c(\Bun_G(\Fq)). \]
由这一描述可知，它把带 Weil 下降数据 $\alpha$ 的紧对象 $\cF\in D_{\Nilp}(\Bun_G)$ 的类
\[ [\cF,\alpha]\in\tr(\Frob|D_{\Nilp}(\Bun_G)) \]
送到 $\cA_c(\Bun_G(\Fq))$ 中预期的函数
\[ x\mapsto\tr(\Frob_x|\cF_x). \]

虽然定理 8.2 是纯粹自守一侧的定理，它的证明却用到谱作用，特别是用到关于 shtuka 空间上同调的结果。

更确切地说，$D_{\Nilp}(\Bun_G)$ 的自对偶数据来自 $D(\Bun_G)$ 的自对偶数据，办法是应用 Beilinson 谱投影子。由于 Beilinson 谱\srcpage{40}投影子由谱作用、从而由 Hecke 作用得到，结果表明只需计算
\[ \tr(\Frob\circ H|D(\Bun_G\times X^{I})) \]
对 Hecke 算子 $H=H_V$，其中 $V\in\Rep(\Gd^{I})$。更确切地说，这必须在 $X^{I}$ 上的 $(\infty,2)$-核范畴中计算，因而结果不是（复形的）向量空间，而是 $D(X^{I})$ 的一个对象。Beilinson 谱投影子与从 $D(X^{I})$ 投影到 $\Ind\Lisse(X)^{\otimes I}$ 有关。所需的关键结果是
\[ \tr(\Frob\circ H|D(\Bun_G\times X^{I}))\in\Ind\Lisse(X)^{\otimes I}\subset D(X^{I}) \]
从而投影操作不起任何作用。

仿照 $H=\id$ 的论证的一个版本，一般地 $\tr(\Frob\circ H|D(\Bun_G\times X^{I}))$ 计算 shtuka 空间的紧支撑上同调。事实上，$H$ 来自 Hecke 对应
\[ \Hck^{I}_{G}\to\Bun_G\times\Bun_G\times X^{I}, \]
而人们必须与 Frobenius 的图像 $\Frob:\Bun_G\to\Bun_G$ 取纤维积。这一纤维积恰是一个 shtuka 空间，它参数化 $X$ 上的 $G$-丛连同它与自己的 Frobenius 拉回之间的一个同构，其中该同构允许在 $X$ 的 $I$ 个点处有极点。借助 Xue（2025）的定理——shtuka 空间的紧支撑上同调是 $X^{I}$ 上的 ind-lisse 层——即可得出结论，且它实际上落在 $\Ind\Lisse(X)^{\otimes I}\subset D(X^{I})$ 中。

\begin{remark}
对定理 1.7，还需要计算谱一侧的 Frobenius 迹。这是直接得多的计算，用到 Gaitsgory 与 Rozenblyum（2017）关于归纳凝聚层发展出的机器。
\end{remark}
